\documentclass[11pt]{article}

\usepackage[margin=1.1in]{geometry}
\usepackage{mathpazo}
\usepackage[T1]{fontenc}
\usepackage[expansion=false]{microtype}
\usepackage{amsmath,amssymb,amsthm,mathtools}
\usepackage{enumitem}
\usepackage{booktabs}
\usepackage[colorlinks=true,linkcolor=black,citecolor=black,urlcolor=black]{hyperref}

\emergencystretch=2.5em

% ---------------------------------------------------------------- theorems
\newtheorem{theorem}{Theorem}[section]
\newtheorem{lemma}[theorem]{Lemma}
\newtheorem{proposition}[theorem]{Proposition}
\newtheorem{corollary}[theorem]{Corollary}
\theoremstyle{definition}
\newtheorem{definition}[theorem]{Definition}
\newtheorem{example}[theorem]{Example}
\newtheorem{hypothesis}[theorem]{Hypothesis}
\theoremstyle{remark}
\newtheorem{remark}[theorem]{Remark}

% ---------------------------------------------------------------- notation
\newcommand{\calc}{\rho_G}
\newcommand{\nil}{\mathbf{0}}
\newcommand{\chaos}{\mathbf{1}}
\newcommand{\rcv}[3]{\mathtt{for}(#1 \leftarrow #2)#3}
\newcommand{\out}[3]{#1!(\langle #2, #3\rangle)}
\newcommand{\qn}[2]{@\langle #1, #2\rangle}
\newcommand{\act}[2]{#1\cdot #2}
\newcommand{\red}{\rightsquigarrow}
\newcommand{\scong}{\equiv}
\newcommand{\FN}{\mathrm{FN}}
\newcommand{\Stab}{\mathrm{Stab}}
\newcommand{\code}{\mathrm{code}}
\newcommand{\Names}{\mathcal{N}}
\newcommand{\match}{\mathrm{match}}
\newcommand{\nf}{\mathrm{nf}}
\newcommand{\Proc}{\mathcal{P}}
\newcommand{\ring}{\mathrm{Ring}}

\title{\textbf{Oriented Processes}\\[0.4em]
\large Group-valued data, a chaos constant, and an equivariant COMM\\
for the reflective higher-order calculus}
\author{L.\,G.\ Meredith\\ \small F1R3FLY.io}
\date{October 18, 2026 \quad\textendash\quad Draft 1}

\begin{document}
\maketitle

\begin{abstract}
\noindent
We take the pattern-matching variant of the rho calculus parametric in a type
of data, and specialise the data to a group $G$. The group acts on processes
freely: an element $g$ multiplies the datum of every name, descending
recursively through quotation. To the calculus we add an annihilating constant
$\chaos$, read as chaos or divergence, with $\chaos \mid P \scong \chaos$,
and an idempotent summation. The communication rule is generalised so that
a receive and a send on two orientations of the same channel produce a sum of
the continuation transported by the relative orientation and by its inverse.
We show that the action is an automorphism of the calculus, that chaos is
forced to have every name free, and that reduction is $G$-equivariant for an
arbitrary group exactly when the two coefficients are chosen as
$\nu(x)\nu(u)^{-1}$ and $\nu(u)\nu(x)^{-1}$. The free action then yields an
orbit-counting theorem: the symmetry group of a process with finitely many
free names acts freely on them, so its order divides their number. For
rotation groups this is a crystallographic restriction on processes. Three
applications indicate the utility of the structure: symmetry reduction for
state-space exploration, where the symmetry is read off the term rather than
declared; symmetric rings, including the dining philosophers, as orbit
products; and \emph{twisted roots} of processes, which generalise the $n$-th
roots of radical fault tolerance and replace its divisibility obstruction by a
symmetry condition in which the stabiliser of a component plays the part of
its replication factor.
\end{abstract}

\tableofcontents

% =====================================================================
\section{Introduction}

Earlier work treated the syntax of the rho calculus as a type constructor
parametric in a type $D$ of data carried by names and messages, and settled on
a pattern-matching binder that lets a receiver bind every component of what it
receives, so that the datum on a message is no longer observationally
inert~\cite{wsh}. That construction asks nothing of $D$. This note asks what
the calculus gains when $D$ is a group.

The answer begins with an action. A group acts on itself by left
multiplication, and that action lifts to processes by multiplying the datum of
every name a process mentions, at every depth of quotation. The lifted action
is free on names: no non-identity element fixes a name. Freeness turns
questions about the symmetry of a process into counting arguments, and the
calculus acquires a notion of \emph{orientation}. The same channel code may
occur at many orientations, and a rotation of a process moves all of them at
once.

Two further ingredients complete the setting. The first is a constant
$\chaos$ for chaos: the process that may do anything, so that composing
anything with it adds nothing. The second is a summation, which the
communication rule needs because its reduct is a sum. The rule itself is the
point of the construction. A receive on one orientation of a channel and a
send on another communicate, and the continuation is delivered twice, once
transported by the relative orientation of the two ends and once by its
inverse. When the two ends agree, the transports cancel and the ordinary rho
communication is recovered.

\paragraph{Contributions.}
\begin{itemize}[leftmargin=1.5em]
\item The calculus $\calc$ (Section~\ref{sec:calculus}): syntax, the free
  action, chaos, idempotent summation, matching, and the oriented COMM rule.
\item Basic observations (Section~\ref{sec:obs}): the action is an
  automorphism; chaos has every name free; matching and substitution commute
  with the action; reduction is $G$-equivariant
  (Theorem~\ref{thm:equivariance}); the calculus is conservative over the
  original; and an orbit-counting theorem
  (Theorem~\ref{thm:counting}) with a crystallographic corollary for rotation
  groups (Corollary~\ref{cor:crystal}).
\item Three applications: symmetry reduction with symmetry computed from the
  term (Section~\ref{sec:reduction}); symmetric rings as orbit products
  (Section~\ref{sec:rings}); and twisted roots of processes, with a
  divisibility theorem that interpolates between replication and sharding
  (Section~\ref{sec:roots}).
\end{itemize}

\paragraph{Scope.}
Summation is taken idempotent throughout, and chaos absorbs it. Other readings
of summation are possible and are not pursued here.

% =====================================================================
\section{The calculus}\label{sec:calculus}

\subsection{Syntax}

Fix a group $(G,\cdot,e,(-)^{-1})$ and a set $V$ of variables, divided into
process variables $X,Y,\dots$ and data variables $\delta,\delta',\dots$. We
write $e$ for the identity of $G$, reserving $\chaos$ for chaos. The terms of
$\calc$ are
\begin{align*}
d &\;::=\; g \;\mid\; \delta
  && g\in G \quad\text{(data)}\\
x,u,\pi &\;::=\; \qn{P}{d}
  && \text{(names; patterns)}\\
P,Q &\;::=\; \nil \;\mid\; \chaos \;\mid\; X \;\mid\; \rcv{\pi}{x}{P}
  \;\mid\; \out{x}{Q}{d} \;\mid\; P \mid Q \;\mid\; P + Q
  && \text{(processes)}
\end{align*}
A name $\qn{P}{d}$ pairs a \emph{code} $P$ with a datum $d$; we write
$\code(\qn{P}{d}) = P$ and $\nu(\qn{P}{d}) = d$. A message
$\out{x}{Q}{d}$ ships the pair $\langle Q, d\rangle$ on $x$; at communication
it is received as the name $\qn{Q}{d}$.

In a receive $\rcv{\pi}{x}{P}$ the pattern $\pi$ is a name in which every
variable occurring is bound by the receive, and each occurs exactly once. The
canonical pattern $\qn{X}{\delta}$ binds both the code and the datum of what
is received. As in~\cite{wsh}, process variables play the part the drop
operator plays in the original calculus~\cite{rho}: a received code is used by
writing its variable in process position. A name is \emph{closed} when it
contains no variables, and a process is closed when every variable in it is
bound. Reduction is defined on closed processes.

\begin{remark}[Notation]
The sketch that motivated this note wrote the action with $*$. Since $*$ is
the drop operator in the rho calculus, we write the action as $\act{g}{P}$.
For the same reason the identity of $G$ is $e$, leaving $\chaos$ free for
chaos.
\end{remark}

\subsection{The free action}

\begin{definition}[Action]\label{def:action}
For $g\in G$ define $\act{g}{-}$ on data, names and processes by
\begin{align*}
\act{g}{h} &= gh \quad (h\in G) &
\act{g}{\delta} &= \delta &
\act{g}{X} &= X\\
\act{g}{\nil} &= \nil &
\act{g}{\chaos} &= \chaos &
\act{g}{\qn{P}{d}} &= \qn{\act{g}{P}}{\act{g}{d}}
\end{align*}
\begin{align*}
\act{g}{(P\mid Q)} &= \act{g}{P}\mid\act{g}{Q} &
\act{g}{(P+Q)} &= \act{g}{P}+\act{g}{Q}\\
\act{g}{\rcv{\pi}{x}{P}} &= \rcv{\act{g}{\pi}}{\act{g}{x}}{\act{g}{P}} &
\act{g}{\out{x}{Q}{d}} &= \out{\act{g}{x}}{\act{g}{Q}}{\act{g}{d}}
\end{align*}
\end{definition}

The action multiplies every literal datum on the left, at every depth of
quotation, inside patterns as well as in names and messages, and leaves
variables alone. Because variables form their own sort, the action never
touches a binding occurrence and needs no side condition against capture.
Equivalently, it is the image of left multiplication
$L_g : G \to G$ under the functor that the parametric syntax defines; the
``free names'' description and the functorial description coincide once the
action descends into codes and patterns.

\subsection{Structural congruence}

Structural congruence $\scong$ is the least congruence, closed under every
term former including quotation, that satisfies the axioms of
Table~\ref{tab:scong}. Two names are equal when their codes are structurally
congruent and their data are equal.

\begin{table}[ht]
\centering
\begin{tabular}{@{}lll@{}}
\toprule
(PAR) & $P\mid\nil \scong P$,\; $P\mid Q\scong Q\mid P$,\;
        $(P\mid Q)\mid R\scong P\mid(Q\mid R)$ & monoid\\
(SUM) & $P+Q\scong Q+P$,\; $(P+Q)+R\scong P+(Q+R)$,\; $P+P\scong P$
      & semilattice\\
(DIST) & $P\mid(Q+R)\scong P\mid Q + P\mid R$ & distributivity\\
(CHAOS) & $\chaos\mid P\scong\chaos$,\; $\chaos + P \scong \chaos$ &
  chaos absorbs both\\
\bottomrule
\end{tabular}
\caption{Axioms of structural congruence.}
\label{tab:scong}
\end{table}

The axiom $\chaos\mid P\scong\chaos$ is the defining property of chaos: it
represents all possible behaviours, so composing it with any $P$ cannot
enrich it. The same reasoning motivates $\chaos+P\scong\chaos$. The two are
consistent with (DIST):
$(\chaos+P)\mid Q\scong\chaos\mid Q + P\mid Q\scong\chaos$. Note that $\nil$ is
not a unit for $+$; it is the inert configuration, one alternative among
others.

By (DIST), every process is congruent to a finite sum of sum-free
\emph{configurations}, or to $\chaos$. Summation is therefore an ensemble of
alternatives, and reduction (below) refines one member of the ensemble.

\subsection{Matching}

Matching is performed on normal forms. Orient (CHAOS) and (DIST) from left to
right and remove duplicate summands; write $\nf(P)$ for a normal form of $P$
modulo associativity and commutativity of $\mid$ and $+$ (AC). A substitution
$\sigma$ maps process variables to closed processes and data variables to
elements of $G$.

\begin{definition}[Matching]\label{def:match}
For a pattern $\pi$ and a closed name $n$,
\[
\match(\pi,n) \;=\; \{\,\sigma \;:\; \pi\sigma =_{AC} \nf(n)\,\}.
\]
\end{definition}

Matching modulo AC can have several solutions, and a communication may use any
of them. One convention in the definition is deliberate.

\begin{remark}[Chaos in quotes]\label{rem:chaosmatch}
Matching is \emph{not} performed modulo (CHAOS). Otherwise the pattern
$\qn{X\mid Y}{\delta}$ would match $\qn{\chaos}{d}$ under every substitution
$\{\chaos/X, Q/Y\}$, since $\chaos\mid Q\scong\chaos$ for all $Q$, and matching
against quoted chaos would not be finitary. In normal form chaos is an atom:
it matches a bare process variable and nothing composite.
\end{remark}

\subsection{Reduction}

The \emph{relative orientation} of a receive on $x$ and a send on $u$ is
$\nu(x)\nu(u)^{-1}$. The reduction relation $\red$ on closed processes is the
least relation closed under the rules of Table~\ref{tab:red}.

\begin{table}[ht]
\centering
\renewcommand{\arraystretch}{1.6}
\begin{tabular}{@{}ll@{}}
\toprule
(COMM) & $\dfrac{\code(x)\scong\code(u) \qquad \sigma\in\match(\pi,\qn{Q}{d})}
  {\rcv{\pi}{x}{P}\mid\out{u}{Q}{d}\;\red\;
   \act{(\nu(x)\nu(u)^{-1})}{P\sigma} \;+\;
   \act{(\nu(u)\nu(x)^{-1})}{P\sigma}}$\\
(PAR) & $\dfrac{P\red P'}{P\mid R\red P'\mid R}$
\qquad (SUM) \; $\dfrac{P\red P'}{P+R\red P'+R}$\\
(STRUCT) & $\dfrac{P\scong P' \qquad P'\red Q' \qquad Q'\scong Q}{P\red Q}$
\qquad (DIV) \; $\chaos\red\chaos$\\
\bottomrule
\end{tabular}
\caption{Reduction rules.}
\label{tab:red}
\end{table}

The premise $\code(x)\scong\code(u)$ is the \emph{gate} of the rule. A receive
and a send interact when they are on the same channel code, possibly at
different orientations, and the relative orientation then acts on the
continuation. The two coefficients are mutually inverse, so writing
$h = \nu(x)\nu(u)^{-1}$ the reduct is $\act{h}{P\sigma}+\act{h^{-1}}{P\sigma}$.

\begin{remark}[The gate]\label{rem:gate}
Other gates are possible.
\begin{itemize}[leftmargin=1.5em]
\item The \emph{strict} gate $x = u$ makes every coefficient $e$, and the
  calculus collapses to the original one extended by chaos and summation.
\item The \emph{open} gate, always true, lets every receive meet every send
  whose message matches its pattern. Pattern matching is then the only
  selector of partners, and channel discipline is lost.
\item A \emph{$K$-gate}, $\code(x)\scong\code(u)$ and
  $\nu(x)\nu(u)^{-1}\in K$ for a subset $K\subseteq G$ that contains $e$,
  is closed under inverses and is closed under conjugation, restricts
  interaction between orientations of a channel to the Cayley graph of $G$
  with connection set $K$.
\end{itemize}
All of the results below hold for each of these gates; we fix the coded gate
as the default because it is the least change that lets orientation matter.
\end{remark}

\begin{remark}[Order of the second coefficient]\label{rem:order}
The sketch that motivated this note used $\nu(x)^{-1}\nu(u)$ as the second
coefficient. In an abelian group it coincides with $\nu(u)\nu(x)^{-1}$. In a
non-abelian group it is the $\nu(x)$-conjugate of the inverse of the first
coefficient, and with that choice reduction is equivariant only for elements
commuting with $\nu(x)^{-1}\nu(u)$; see Remark~\ref{rem:nonabelian}. The choice
$\nu(u)\nu(x)^{-1}$ is the one that makes reduction equivariant for every
group, and in particular for non-abelian rotation groups such as $SO(3)$.
\end{remark}

% =====================================================================
\section{Basic observations}\label{sec:obs}

\subsection{The action is an automorphism}

\begin{lemma}\label{lem:auto}
\begin{enumerate}[label=(\roman*)]
\item $\act{e}{P}=P$ and $\act{g}{(\act{h}{P})}=\act{(gh)}{P}$.
\item $P\scong Q$ if and only if $\act{g}{P}\scong\act{g}{Q}$.
\item $\nf(\act{g}{P}) =_{AC} \act{g}{\nf(P)}$.
\end{enumerate}
So each $\act{g}{-}$ is an automorphism of
$(\Proc/{\scong},\mid,+,\nil,\chaos)$, and $G$ acts on processes up to
structural congruence.
\end{lemma}

\begin{proof}
(i) is by induction on terms, the only base case being a literal datum, where
$g(hd)=(gh)d$. For (ii), the action maps every instance of an axiom in
Table~\ref{tab:scong} to an instance of the same axiom, because it commutes
with the term formers and fixes $\nil$ and $\chaos$. Hence it preserves
$\scong$, and $\act{g^{-1}}{-}$ is its inverse by (i). (iii) holds because the
oriented normalisation rules are themselves invariant under the action.
\end{proof}

\subsection{Chaos has every name}

Free names collect the closed names that occur in subject position or in a
payload. Writing $\Names$ for the set of closed names,
\begin{align*}
\FN(\nil)&=\FN(X)=\emptyset, &
\FN(P\mid Q)&=\FN(P+Q)=\FN(P)\cup\FN(Q),\\
\FN(\rcv{\pi}{x}{P})&=(\{x\}\cap\Names)\cup\FN(P), &
\FN(\out{x}{Q}{d})&=(\{x\}\cap\Names)\cup\FN(Q),
\end{align*}
and $\FN(\chaos)$ is fixed by the following proposition.

\begin{proposition}\label{prop:chaosnames}
If $\FN$ respects $\scong$ and is a union over $\mid$, then
$\FN(\chaos)=\Names$.
\end{proposition}

\begin{proof}
For any closed name $x$, $\FN(\chaos)=\FN(\chaos\mid\out{x}{\nil}{e})
=\FN(\chaos)\cup\{x\}$, so $x\in\FN(\chaos)$.
\end{proof}

Chaos has every name free because it is absorbing. Two consequences are worth
recording at once. First, $\act{g}{\chaos}=\chaos$ is no longer an axiom about
the action but an instance of its definition, since left multiplication
permutes $\Names$. Second, no name is fresh for chaos: an argument that picks
a channel not free in $P$ fails as soon as $P\scong\chaos$.

\begin{lemma}\label{lem:fnaction}
$\FN(\act{g}{P})=\act{g}{\FN(P)}$, and $\FN$ respects $\scong$.
\end{lemma}

\begin{proof}
The first claim is by induction: the action preserves closedness of names and
permutes $\Names$. The second is a check of the axioms; for (CHAOS),
$\FN(\chaos\mid P)=\Names\cup\FN(P)=\Names=\FN(\chaos)$.
\end{proof}

\subsection{Matching and substitution commute with the action}

For a substitution $\sigma$ write $\act{g}{\sigma}$ for the substitution
$X\mapsto\act{g}{\sigma(X)}$, $\delta\mapsto g\,\sigma(\delta)$.

\begin{lemma}\label{lem:subst}
$\act{g}{(P\sigma)}=(\act{g}{P})(\act{g}{\sigma})$, and
$\sigma\in\match(\pi,n)$ if and only if
$\act{g}{\sigma}\in\match(\act{g}{\pi},\act{g}{n})$.
\end{lemma}

\begin{proof}
The substitution identity is by induction on $P$. At a process variable both
sides give $\act{g}{\sigma(X)}$. At a data variable in data position, the
action fixes $\delta$ on the right, which the substitution then sends to
$g\,\sigma(\delta)$, as on the left. At a literal datum both sides give $gh$.
For matching, $\pi\sigma=_{AC}\nf(n)$ if and only if
$\act{g}{(\pi\sigma)}=_{AC}\act{g}{\nf(n)}$, because the action is a bijection
that respects AC; by the substitution identity and
Lemma~\ref{lem:auto}(iii) this reads
$(\act{g}{\pi})(\act{g}{\sigma})=_{AC}\nf(\act{g}{n})$.
\end{proof}

For example, a pattern with a literal datum, $\qn{X}{h}$, matches
$\qn{Q}{d}$ exactly when $h=d$; after the action it is $\qn{X}{gh}$, which
matches $\qn{\act{g}{Q}}{gd}$ exactly when $gh=gd$. Acting on patterns is what
keeps a rotated redex a redex.

\subsection{Equivariance}

\begin{theorem}[Equivariance]\label{thm:equivariance}
If $P\red P'$ then $\act{g}{P}\red\act{g}{P'}$, for every $g\in G$.
\end{theorem}

\begin{proof}
By induction on the derivation. (PAR), (SUM) and (STRUCT) follow from
Lemma~\ref{lem:auto}, and (DIV) from $\act{g}{\chaos}=\chaos$. For (COMM), let
$a=\nu(x)$ and $b=\nu(u)$. The rotated redex is
$\rcv{\act{g}{\pi}}{\act{g}{x}}{\act{g}{P}}\mid\out{\act{g}{u}}{\act{g}{Q}}{gd}$.
Its gate holds because $\code(\act{g}{x})=\act{g}{\code(x)}$ and
Lemma~\ref{lem:auto}(ii). By Lemma~\ref{lem:subst}, $\act{g}{\sigma}$ is a
matcher, and the continuation becomes
$(\act{g}{P})(\act{g}{\sigma})=\act{g}{(P\sigma)}$. The rotated data are $ga$
and $gb$, so the coefficients are
\[
(ga)(gb)^{-1}=g\,ab^{-1}g^{-1},
\qquad
(gb)(ga)^{-1}=g\,ba^{-1}g^{-1}.
\]
By Lemma~\ref{lem:auto}(i) the reduct is
\[
\act{(g\,ab^{-1}g^{-1})}{\act{g}{(P\sigma)}}
+\act{(g\,ba^{-1}g^{-1})}{\act{g}{(P\sigma)}}
=\act{g}{\bigl(\act{(ab^{-1})}{P\sigma}+\act{(ba^{-1})}{P\sigma}\bigr)},
\]
which is $g$ applied to the original reduct.
\end{proof}

\begin{remark}[The non-abelian case]\label{rem:nonabelian}
With $a^{-1}b$ as the second coefficient, the rotated coefficient is
$(ga)^{-1}(gb)=a^{-1}b$, unchanged, and the second summand becomes
$\act{(a^{-1}b\,g)}{P\sigma}$ where equivariance needs
$\act{(g\,a^{-1}b)}{P\sigma}$. These agree for all $P$ only when $g$ commutes
with $a^{-1}b$. Equivariance is therefore what selects the coefficients.
\end{remark}

\subsection{Conservativity and frames}

\begin{proposition}[Conservativity]\label{prop:conservative}
If $\nu(x)=\nu(u)$, in particular if $x=u$, then
$\rcv{\pi}{x}{P}\mid\out{u}{Q}{d}\red P\sigma$. Hence for the trivial group,
and on any process whose communications occur between equal names, $\calc$
reduces exactly as the pattern-matching calculus of~\cite{wsh}, extended by
chaos and summation.
\end{proposition}

\begin{proof}
Both coefficients are $e$, so the reduct is $P\sigma+P\sigma\scong P\sigma$ by
idempotence.
\end{proof}

This is the cancellation in the rule. The relative orientation of matched ends
is the identity, and the doubling of the continuation disappears into the
idempotence of summation.

\paragraph{Symmetry between the ends.}
Exchanging the roles of $\nu(x)$ and $\nu(u)$ exchanges the two summands of the
reduct, which is unchanged. The rule does not depend on whose frame the
relative orientation is measured in.

\paragraph{Frames.}
Consider a name $n$ in the receiver's frame, $\nu(n)=\nu(x)$. The second summand
moves it to $\nu(u)\nu(x)^{-1}\nu(x)=\nu(u)$: into the sender's frame. The first
moves it to $\nu(x)\nu(u)^{-1}\nu(x)$, the reflection of the sender's frame
through the receiver's; in an abelian group written additively, a receiver at
$j$ hearing from $k$ lands at $k$ in one summand and at $2j-k$ in the other.

\paragraph{Anisotropic and isotropic codes.}
A code $c$ is \emph{isotropic} for $H\le G$ when $\act{h}{c}\scong c$ for all
$h\in H$; nameless codes such as $\nil$ are isotropic for all of $G$. All
orientations $\qn{c}{h}$ of an isotropic code share a code, so under the coded
gate they communicate with one another, with non-trivial coefficients. A code
with names is \emph{anisotropic}: rotating a name moves its code as well as
its datum, and by Theorem~\ref{thm:counting} below only finitely many rotations
can fix it. Anisotropic codes therefore produce genuinely distinct channels
under rotation, and processes built from them reduce exactly as in the
original calculus (Proposition~\ref{prop:conservative}). Rotation can thus be
used either to \emph{place} copies of a process on distinct channels, or to
\emph{orient} a single channel.

\subsection{Stabilisers and orbit counting}

\begin{definition}
The \emph{symmetry group} of $P$ is
$\Stab(P)=\{\,g\in G : \act{g}{P}\scong P\,\}$.
\end{definition}

By Lemma~\ref{lem:auto}, $\Stab(P)$ is a subgroup,
$\Stab(\act{g}{P})=g\,\Stab(P)\,g^{-1}$,
$\Stab(P\mid Q)\supseteq\Stab(P)\cap\Stab(Q)$, and
$\Stab(\nil)=\Stab(\chaos)=G$.

\begin{theorem}[Orbit counting]\label{thm:counting}
$\Stab(P)$ acts freely on $\FN(P)$. Consequently, if $\FN(P)$ is finite, then
either $\FN(P)=\emptyset$ and $\Stab(P)=G$, or $\Stab(P)$ is finite and
$|\Stab(P)|$ divides $|\FN(P)|$.
\end{theorem}

\begin{proof}
If $g\in\Stab(P)$, Lemma~\ref{lem:fnaction} gives
$\act{g}{\FN(P)}=\FN(\act{g}{P})=\FN(P)$, so $\Stab(P)$ acts on $\FN(P)$. The
action is free: $\act{g}{\qn{Q}{d}}=\qn{Q}{d}$ requires $gd=d$, hence $g=e$.
If $\FN(P)$ is finite and non-empty, every orbit of $\Stab(P)$ on it has
$|\Stab(P)|$ elements, so $\Stab(P)$ is finite and its order divides
$|\FN(P)|$. If $\FN(P)=\emptyset$, then $P$ contains no receive, no send and no
chaos outside quotation, so it is built from $\nil$ and the two operators and
every $g$ fixes it.
\end{proof}

$\FN(P)$ is finite exactly when $\chaos$ does not occur in $P$ outside quotation.
So the theorem covers every process that is not chaotic, and chaos is the one
way for a process with names to be fixed by an infinite group.

\begin{corollary}\label{cor:infinite}
If $G$ is infinite and $\Stab(P)=G$, then $\FN(P)=\emptyset$ or
$\FN(P)=\Names$. In words: the only processes invariant under an infinite
group are those that mention no names and those that contain chaos.
\end{corollary}

\begin{proposition}[Symmetry is computable]\label{prop:computable}
If $\FN(P)$ is finite and non-empty, fix $n_0\in\FN(P)$. Then
\[
\Stab(P)\subseteq\{\,\nu(n)\,\nu(n_0)^{-1} : n\in\FN(P)\,\},
\]
so $\Stab(P)$ is found by at most $|\FN(P)|$ congruence checks, for any group,
finite or not.
\end{proposition}

\begin{proof}
If $g\in\Stab(P)$ then $\act{g}{n_0}\in\FN(P)$, and
$\nu(\act{g}{n_0})=g\,\nu(n_0)$ determines $g$.
\end{proof}

\subsection{Rotations}

Take the data to be orientations. With $G=SO(2)$ a name is a channel code at a
planar orientation, and with $G=SO(3)$ at an orientation in space.

\begin{corollary}[Crystallographic restriction]\label{cor:crystal}
Let $\FN(P)$ be finite and non-empty.
\begin{enumerate}[label=(\roman*)]
\item If $G=SO(2)$, then $\Stab(P)$ is cyclic of order dividing $|\FN(P)|$.
\item If $G=SO(3)$, then $\Stab(P)$ is cyclic, dihedral, tetrahedral,
  octahedral or icosahedral, of order dividing $|\FN(P)|$.
\end{enumerate}
\end{corollary}

\begin{proof}
By Theorem~\ref{thm:counting} the stabiliser is a finite subgroup, and these
are the finite subgroups of $SO(2)$ and $SO(3)$~\cite{armstrong}.
\end{proof}

A process can have the rotational symmetry of a polygon or a polyhedron, but
never continuous rotational symmetry unless it mentions no names or is
chaotic. An icosahedrally symmetric process mentions a multiple of $60$ names.

\subsection{Orbit products and orbit sums}

For a finite subgroup $H\le G$ and a process $R$ define
\[
\textstyle\prod_H R \;=\; \big|_{h\in H}\;\act{h}{R},
\qquad
\sum_H R \;=\; \sum_{h\in H}\;\act{h}{R}.
\]
Both are $H$-invariant: $\act{k}{\prod_H R}$ re-indexes the product by
$h\mapsto kh$, and associativity and commutativity of $\mid$ do the rest; the
same holds for $\sum_H R$. Orbit products are the standard way to impose a
symmetry on a process, with $R$ playing the part of a fundamental domain. The
reduct of (COMM) is a partial orbit sum, over the two-element set
$\{h,h^{-1}\}$, and is invariant under the inversion that exchanges them.

% =====================================================================
\section{Application I: symmetry reduction}\label{sec:reduction}

Explicit-state exploration of concurrent systems is limited by the size of the
state space, and symmetry reduction is one of the standard
remedies~\cite{emerson-sistla,ip-dill,clarke-sym}: when a group acts on the
states by automorphisms of the transition relation, it suffices to explore one
representative per orbit. In those treatments the symmetry is a property of
a particular system and must be declared or detected. In $\calc$ it is a
property of the calculus.

\begin{theorem}[Orbit quotient]\label{thm:quotient}
Let $H\le G$. Define $[P]=\{\act{h}{P'} : h\in H,\, P'\scong P\}$ and
$[P]\red_H[Q]$ if $P\red Q'$ for some $Q'\in[Q]$. Then $\red_H$ is
well-defined, and the quotient map $P\mapsto[P]$ is a functional
bisimulation from $(\Proc,\red)$ to $(\Proc/H,\red_H)$. Consequently, for any
set of states $\Phi$ closed under $\scong$ and under the action of $H$, a state
in $\Phi$ is reachable from $P_0$ if and only if an orbit in $\Phi/H$ is
reachable from $[P_0]$.
\end{theorem}

\begin{proof}
Well-definedness: if $\act{h}{P}\red Q$ then $P\red\act{h^{-1}}{Q}$ by
Theorem~\ref{thm:equivariance}, and $\act{h^{-1}}{Q}\in[Q]$. The bisimulation
clauses are then immediate: every step of $P$ projects to a step of $[P]$, and
every step of $[P]$ is the projection of a step of $P$ by the same argument.
\end{proof}

Three features distinguish this from the usual setting.

\begin{enumerate}[leftmargin=1.5em]
\item \emph{Soundness for every subgroup.} Equivariance holds for the whole
  calculus, not for a given system, so the quotient is sound for any $H$ and
  any initial process. Barbs move under the action, $\act{h}{P}\downarrow
  \act{h}{x}$ whenever $P\downarrow x$, so the invariant observations are
  orbit barbs $P\downarrow H{\cdot}x$, and for a spatial logic whose formulae mention names, with the action extended
  to formulae, satisfaction is equivariant: $\act{g}{P}\models\act{g}{\varphi}$
  whenever $P\models\varphi$. Invariant formulae can be checked on the quotient.
\item \emph{Effectiveness from the term.} The quotient pays when the
  reachable set is closed under $H$, which holds when $H\le\Stab(P_0)$: then
  $P_0\red^*P$ gives $P_0\scong\act{h}{P_0}\red^*\act{h}{P}$. By
  Proposition~\ref{prop:computable} the largest such $H$ is computed from
  $P_0$ with $|\FN(P_0)|$ checks; nothing is declared.
\item \emph{A bound.} By Theorem~\ref{thm:counting}, $|\Stab(P_0)|$ divides
  $|\FN(P_0)|$. The symmetry that can be extracted from the initial state is
  bounded by its number of free names, and the reduction factor on any orbit
  is at most $|H|$.
\end{enumerate}

Canonicalisation, the choice of a representative of $[P]$, costs at most
$|H|$ congruence-class comparisons for finite $H$, by taking the least of
$\{\act{h}{P}\}_{h\in H}$ in any total order on normal forms. A spatial logic
model checker for the rho calculus, such as the port of SLMC~\cite{slmc}, is
the natural home for this reduction, since spatial formulae already range over
the structure that the action permutes.

% =====================================================================
\section{Application II: symmetric rings}\label{sec:rings}

Rings of identical agents are the classical setting for symmetry in
distributed computing. Orbit products give them a uniform syntax.

\subsection{The construction}

Let $r\in G$ have order $n$: a generator of $\mathbb{Z}/n$, or the rotation by
$2\pi/n$ in $SO(2)$. Choose two anisotropic codes, for instance
\[
a_f=\out{\qn{\nil}{e}}{\nil}{e},
\qquad
a_t=\rcv{\qn{X}{\delta}}{\qn{\nil}{e}}{\nil},
\]
and base names $f=\qn{a_f}{e}$ (a fork) and $t=\qn{a_t}{e}$ (a table, on which
eating is announced). Write $f_k=\act{r^k}{f}$ and $t_k=\act{r^k}{t}$; indices
are taken mod $n$. Since $\act{r^k}{a_f}=\out{\qn{\nil}{r^k}}{\nil}{r^k}$, the
codes of $f_0,\dots,f_{n-1}$ are pairwise distinct, so the forks are $n$
distinct channels and, by Proposition~\ref{prop:conservative}, every
communication below is an ordinary one. Let
\begin{align*}
\mathit{Phil} &= \rcv{\qn{X}{\delta}}{f_0}{\rcv{\qn{Y}{\delta'}}{f_1}{\bigl(
  \out{t_0}{\nil}{e}\mid\out{f_0}{X}{\delta}\mid\out{f_1}{Y}{\delta'}\bigr)}},\\
\mathit{Fork} &= \out{f_0}{\nil}{e},
\qquad
R=\mathit{Phil}\mid\mathit{Fork},
\qquad
\ring_n=\textstyle\prod_{\langle r\rangle}R.
\end{align*}
The copy $\act{r^k}{R}$ is the $k$-th philosopher, who takes fork $f_k$ and
then fork $f_{k+1}$, announces on $t_k$, and returns both forks, together with
fork $f_k$ itself, carrying the token $\langle\nil,r^k\rangle$. The pattern
$\qn{X}{\delta}$ is fixed by the action, so every philosopher accepts every
token.

\begin{proposition}\label{prop:ring}
\begin{enumerate}[label=(\roman*)]
\item $\Stab(\ring_n)=\langle r\rangle\cong\mathbb{Z}/n$ when $G=\mathbb{Z}/n$
  or $G=SO(2)$.
\item $\ring_n$ reduces in $n$ steps to the deadlocked state
  $D_n=\prod_{\langle r\rangle}W$, where
  $W=\rcv{\qn{Y}{\delta'}}{f_1}{\bigl(\out{t_0}{\nil}{e}\mid
  \out{f_0}{\nil}{e}\mid\out{f_1}{Y}{\delta'}\bigr)}$.
  The state $D_n$ has no reduct, has no barb on any $t_k$, and
  $\Stab(D_n)=\langle r\rangle$.
\item For $n\ge2$, $\ring_n$ also reaches states with a barb on some $t_k$.
\end{enumerate}
\end{proposition}

\begin{proof}
(i) The product is $\langle r\rangle$-invariant. Conversely
$\FN(\ring_n)=\{f_k,t_k\}_{k}$ has $2n$ elements. A symmetry maps fork names
to fork names, since a send code is never congruent to a receive code, so
by Proposition~\ref{prop:computable} it is $\nu(f_k)\nu(f_0)^{-1}=r^k$.
(ii) For each $k$ in turn, philosopher $k$ receives the token of fork $k$; the
coefficients are $e$, $\sigma=\{\nil/X,\,r^k/\delta\}$, and the continuation is
$\act{r^k}{W}$. After $n$ such steps no send on any $f_k$ remains, so no
(COMM) applies. Its symmetry group is computed as in (i). (iii) Let
philosopher $0$ take $f_0$ and then $f_1$ before philosopher $1$ moves.
\end{proof}

The $n$ steps to $D_n$ are a single step and its $n-1$ rotations: the $k$-th is
$\act{r^k}{}$ of the $0$-th, placed in context. Equivariance is exactly the
property that the adversary of Lehmann and Rabin exploits~\cite{lehmann-rabin}:
from a symmetric state, a scheduler that lets every copy take the
corresponding step returns the system to a symmetric state, and a deterministic
symmetric protocol cannot leave the symmetric states without help. Their
theorem, that no deterministic, fully distributed and symmetric solution to the
dining philosophers exists, can therefore be stated inside $\calc$ with
$\langle r\rangle$-invariance as the formal notion of symmetry. Its proof in
the calculus is left to further work, together with the question of which
ingredient of $\calc$ is the symmetry breaker that their randomised solution
requires. With summation read as an idempotent ensemble, as here,
$\mathit{Phil}_{L}+\mathit{Phil}_{R}$ is demonic choice rather than a coin, and
the deadlocked configuration remains in the ensemble.

\begin{remark}[Isotropic rings]
If the anisotropic code $a_f$ is replaced by the isotropic code $\nil$, the
names $\qn{\nil}{r^k}$ all share a code. Under the coded gate every orientation
then communicates with every other, and a receive at $r^j$ from a send at
$r^k$ transports the names of its own frame to $r^k$ in one summand and to
$r^{2j-k}$ in the other. Under the $K$-gate with $K=\{e,r,r^{-1}\}$ only neighbours
communicate, and the interaction topology is the Cayley graph of
$\mathbb{Z}/n$. Rings of isotropic channels are therefore not rings of fixed
agents but of agents that move as they communicate.
\end{remark}

% =====================================================================
\section{Application III: twisted roots}\label{sec:roots}

The note on radical fault tolerance~\cite{rft} studies $n$-th roots of
processes, $R$ with $R^n\approx P$, where $R^n$ is the $n$-fold parallel
composition. Under a unique-decomposition hypothesis the existence of roots is
governed by divisibility: every prime component of $P$ must occur with
multiplicity divisible by $n$, so that a component occurring once has no
root. The group action offers a second way to divide a process into $n$ equal
parts: into $n$ copies of one process at different orientations.

\begin{definition}
Let $H\le G$ be finite. A process $R$ is an \emph{$H$-twisted root} of $P$ if
$\prod_H R\scong P$.
\end{definition}

If $H\le\Stab(R)$ the twist is trivial, $\prod_H R\scong R^{|H|}$, and an
untwisted root is recovered. At the other extreme, if $f_k=\act{r^k}{f}$ as in
Section~\ref{sec:rings}, the process
\[
P=\out{f_0}{\nil}{e}\mid\out{f_1}{\nil}{r}\mid\cdots\mid
\out{f_{n-1}}{\nil}{r^{n-1}}
\]
has every prime component with multiplicity one and so no $n$-th root, yet
$\out{f_0}{\nil}{e}$ is a $\mathbb{Z}/n$-twisted root of it. To say exactly when
twisted roots exist we adopt the hypothesis of~\cite{rft}.

\begin{hypothesis}[Unique decomposition]\label{hyp:upd}
Every closed process that is free of chaos and summation is structurally
congruent to a parallel composition of primes, unique up to $\scong$ and
permutation, where a prime is a process not congruent to $\nil$ and not
congruent to a composition of two processes neither of which is $\nil$.
\end{hypothesis}

Unique decomposition is known for several process calculi~\cite{milner-moller};
for the rho calculus it is assumed, as in~\cite{rft}. Write $M_P$ for the
multiset of primes of $P$. By Lemma~\ref{lem:auto} the action maps primes to
primes, so $H$ acts on primes and on such multisets, and each prime $q$ has a
stabiliser $\Stab_H(q)=\Stab(q)\cap H$.

\begin{theorem}[Twisted divisibility]\label{thm:twisted}
Assume Hypothesis~\ref{hyp:upd}, and let $H\le G$ be finite and $P$ free of
chaos and summation. Then $P$ has an $H$-twisted root if and only if
\begin{enumerate}[label=(\alph*)]
\item $M_P$ is $H$-invariant, and
\item for every prime $q$, $|\Stab_H(q)|$ divides the multiplicity $M_P(q)$.
\end{enumerate}
\end{theorem}

\begin{proof}
($\Rightarrow$) If $\prod_H R\scong P$ then, by uniqueness,
$M_P=\sum_{h\in H}\act{h}{M_R}$, which is invariant. For a prime $q$,
\[
M_P(q)=\sum_{h\in H}M_R(\act{h^{-1}}{q})
=|\Stab_H(q)|\sum_{q'\in H\cdot q}M_R(q'),
\]
because each $q'$ in the orbit of $q$ equals $\act{h^{-1}}{q}$ for exactly
$|\Stab_H(q)|$ elements $h$.
($\Leftarrow$) Choose one prime $q$ from each $H$-orbit in $M_P$ and let $R$
be the composition of the chosen primes, each with multiplicity
$M_P(q)/|\Stab_H(q)|$. By the same count, $\prod_H R$ gives each $q'$ in the
orbit of a chosen $q$ multiplicity
$|\Stab_H(q')|\cdot M_P(q)/|\Stab_H(q)|=M_P(q)=M_P(q')$, using that stabilisers
along an orbit are conjugate and that $M_P$ is invariant.
\end{proof}

\begin{corollary}\label{cor:roots}
Under the hypotheses of Theorem~\ref{thm:twisted}:
\begin{enumerate}[label=(\roman*)]
\item If $H$ fixes every prime of $P$, then $P$ has an $H$-twisted root if and
  only if $|H|$ divides every multiplicity. This is the root criterion
  of~\cite{rft} with $n=|H|$.
\item If $H$ acts freely on the primes of $P$, then $P$ has an $H$-twisted root
  if and only if $M_P$ is $H$-invariant: symmetry replaces divisibility.
\item By Theorem~\ref{thm:counting}, $|\Stab_H(q)|$ divides
  $\gcd(|H|,|\FN(q)|)$. A prime whose number of free names is coprime to $|H|$
  imposes no divisibility condition.
\item Chaos is its own twisted root for every $H$: $\prod_H\chaos\scong\chaos$.
\end{enumerate}
\end{corollary}

\subsection{The stabiliser is the replication factor}

Theorem~\ref{thm:twisted} interpolates between two familiar ways of spreading
a system over $n$ sites. Take an orbit $O$ of primes with stabiliser $S$. In
$\prod_H R$, each prime of $O$ is contributed by exactly $|S|$ of the copies.
\begin{itemize}[leftmargin=1.5em]
\item When $S=H$ every copy contributes the same component: this is
  \emph{replication}, and the copies are those of an ordinary root.
\item When $S$ is trivial each component comes from exactly one copy, and the
  copies partition the system among themselves: this is \emph{sharding}, in
  which every site runs the same program in its own frame.
\end{itemize}
In general the replication factor of a component is the order of its
stabiliser, and by Corollary~\ref{cor:roots}(iii) it divides both $|H|$ and its
number of free names.

This bears directly on the arbitration problem of~\cite{rft}. There, copies of
a root contend for the same linear messages, and an arbiter must be installed
at every contention site, at a cost of $\Theta(n)$ rendezvous per site. In a
twisted root, identical primes, and hence identical channels, are shared only
among the $|S|$ copies that replicate them. The contention that requires
arbitration is confined to orbits with non-trivial stabiliser, and its width
is $|S|$ rather than $|H|$.

\paragraph{Coupling.}
Copies with trivial stabilisers can still interact, through channels they
share. Let
\[
C(R)=\{\,h\in H\setminus\{e\} : \code(m)\scong\act{h}{\code(n)}
\text{ for some } m,n\in\FN(R)\,\}.
\]
In the initial configuration, copies $\act{k}{R}$ and $\act{k'}{R}$ have a
redex between them only if $k^{-1}k'\in C(R)$, so the static coupling between
copies is the Cayley graph of $H$ with connection set $C(R)$. For the ring of
Section~\ref{sec:rings}, $C(R)=\{r,r^{-1}\}$, since philosopher $k$ shares fork
$f_{k+1}$ with philosopher $k+1$; $R$ is a twisted root of $\ring_n$ with trivial
stabilisers, and its coupling graph is the cycle. When $C(R)=\emptyset$
the copies are disjoint in the initial configuration. Whether they stay
disjoint depends on the names they can construct at run time, which is the
namespace question that rho always raises and that the action does not settle.

\paragraph{Rotations again.}
With $G=SO(3)$ and $H$ the icosahedral group, a twisted root spreads a system
over $60$ sites at the orientations of the icosahedron. By
Corollary~\ref{cor:roots}(iii), a component with a number of free names
coprime to $60$ is sharded, never replicated.

% =====================================================================
\section{Further work}\label{sec:further}

\begin{itemize}[leftmargin=1.5em]
\item \emph{Behavioural twisted roots.} The structural theorem above should be
  lifted to $\prod_H R\approx P$ for an equivariant behavioural equivalence,
  relative to an observer as in~\cite{rft}, with orbit barbs as observations.
\item \emph{Arbitration in twisted roots.} The claim that arbitration is
  confined to orbits with non-trivial stabiliser should be made precise by
  combining the arbiter of~\cite{rft} with the coupling graph.
\item \emph{Lehmann--Rabin in $\calc$.} State and prove the impossibility of
  deterministic symmetric solutions with $\langle r\rangle$-invariance as the
  notion of symmetry, and identify the symmetry breaker that a randomised
  solution needs. A probabilistic reading of summation is one candidate.
\item \emph{Transport on isotropic channels.} The behaviour of isotropic
  rings under the coded gate and under $K$-gates, where communication moves
  agents along a Cayley graph.
\item \emph{Symmetry reduction in practice.} Implement stabiliser computation
  and orbit canonicalisation in the spatial logic model checker, and measure
  the reduction on rings and twisted roots.
\item \emph{Unique decomposition.} Hypothesis~\ref{hyp:upd} is inherited
  from~\cite{rft} and remains open for the rho calculus.
\end{itemize}

\section*{Acknowledgements}
This note was drafted in collaboration with Claude (Anthropic).

\begin{thebibliography}{99}
\bibitem{armstrong}
M.~A. Armstrong.
\newblock \emph{Groups and Symmetry}.
\newblock Undergraduate Texts in Mathematics. Springer, 1988.

\bibitem{slmc}
L.~Caires and L.~Cardelli.
\newblock A spatial logic for concurrency (part {I}).
\newblock \emph{Information and Computation}, 186(2):194--235, 2003.

\bibitem{clarke-sym}
E.~M. Clarke, E.~A. Emerson, S.~Jha, and A.~P. Sistla.
\newblock Symmetry reductions in model checking.
\newblock In \emph{CAV}, LNCS 1427, pages 147--158. Springer, 1998.

\bibitem{emerson-sistla}
E.~A. Emerson and A.~P. Sistla.
\newblock Symmetry and model checking.
\newblock \emph{Formal Methods in System Design}, 9(1/2):105--131, 1996.

\bibitem{ip-dill}
C.~N. Ip and D.~L. Dill.
\newblock Better verification through symmetry.
\newblock \emph{Formal Methods in System Design}, 9(1/2):41--75, 1996.

\bibitem{lehmann-rabin}
D.~Lehmann and M.~O. Rabin.
\newblock On the advantages of free choice: a symmetric and fully distributed
  solution to the dining philosophers problem.
\newblock In \emph{POPL}, pages 133--138. ACM, 1981.

\bibitem{rho}
L.~G. Meredith and M.~Radestock.
\newblock A reflective higher-order calculus.
\newblock \emph{Electronic Notes in Theoretical Computer Science},
  141(5):49--67, 2005.

\bibitem{rft}
L.~G. Meredith.
\newblock Radical fault tolerance: roots of processes, arbitration, and the
  resolver that consensus supplies.
\newblock F1R3FLY.io research note, 2026.

\bibitem{wsh}
L.~G. Meredith.
\newblock Where substitution happens.
\newblock F1R3FLY.io research note, 2026.

\bibitem{milner-moller}
R.~Milner and F.~Moller.
\newblock Unique decomposition of processes.
\newblock \emph{Theoretical Computer Science}, 107(2):357--363, 1993.
\end{thebibliography}

\end{document}
